% Book 4, Year 2, chapter 31 -- project: a key-value server in Rust.
% Provenance (never printed): sources/book4/BRIEFS.md, chapter 31.
\chapter{Project: A Key-Value Server}\label{ch:b2:project-key-value-server}

The server of this chapter listens on \code{127.0.0.1:44310} and keeps a map from string keys to string values. It
speaks a line-based text protocol: a client sends \code{SET k v}, \code{GET k}, \code{DEL k}, \code{COUNT} or
\code{QUIT}, one per line, and receives one line for each, \code{OK}, \code{VALUE v}, \code{NOTFOUND}, \code{COUNT n} or
\code{ERR} followed by a reason. It serves many clients at once, it survives a restart by replaying an \omterm{def:b2:project-key-value-server:log}{append-only log},
and it is measured under a \omterm{def:b2:project-key-value-server:load-test}{load test}. It is written in Rust with the standard library only, in six milestones, each with
its code, its output and its tests; the Year 3 volume will make the same server survive the loss of a machine.

\begin{recall}
\omterm{def:b2:rust-ownership:ownership}{Ownership}, borrowing and lifetimes (\cref{ch:b2:rust-ownership}); \code{Result} and the \code{?} operator
(\cref{ch:b2:rust-enums-and-errors}); \omterm{def:b2:rust-traits-and-generics:trait}{traits} and \omterm{def:b2:rust-traits-and-generics:trait-object}{trait objects} (\cref{ch:b2:rust-traits-and-generics}). \omterm{def:b2:threads-and-synchronisation:thread}{Threads},
mutexes, reader--writer \omterm{def:b2:threads-and-synchronisation:lock}{locks}, and the lock-order rule that prevents deadlock
(\cref{prop:b2:threads-and-synchronisation:order}). A \omterm{def:b2:transport-layer:tcp}{TCP connection} carries a stream of bytes with no message boundaries,
so a protocol must frame its messages (\cref{sec:b2:sockets:framing}, \cref{sec:b2:transport-layer:tcp}). The CRC-32
detects every \omterm{def:b2:link-layer:burst}{burst} of errors up to 32 bits long (\cref{thm:b2:link-layer:crc-burst}). The Year 1 volume: hash tables,
whose operations take expected constant time.
\end{recall}

\begin{definition}[Key--value store]\label{def:b2:project-key-value-server:kv}
A \emph{key--value store}\index{key--value store} is a service that keeps a finite map from keys to values and offers
three operations on it: set the value of a key, get the value of a key, and delete a key. It offers no query over the
values, no join and no schema: a client that needs them builds them from keys.
\end{definition}

\begin{omfigure}
\begin{tikzpicture}[b/.style={draw, rounded corners=2pt, font=\small, align=center, minimum height=8mm},
  >={Stealth[length=1.8mm]}]
\foreach \i in {0, 1, 2} \node[b, minimum width=14mm] (c\i) at (0, -\i*1.05) {client \the\numexpr\i+1\relax};
\node[b, minimum width=22mm] (l) at (3.2, -1.05) {listener\\\scriptsize\code{:44310}};
\node[b, minimum width=16mm] (q) at (6.4, -1.05) {queue};
\foreach \i in {0, 1, 2, 3} \node[b, minimum width=16mm, fill=omDef!8] (w\i) at (8.6, 0.3-\i*0.9) {worker \i};
\foreach \i/\y in {0/0.3, 1/-0.6, 2/-1.5} \node[b, minimum width=12mm, fill=omProp!10, font=\scriptsize] (s\i)
  at (11.2, \y) {shard \i};
\node[font=\small] at (11.2, -2.15) {$\vdots$};
\node[b, minimum width=12mm, fill=omProp!10, font=\scriptsize] (s3) at (11.2, -2.9) {shard 15};
\node[b, minimum width=12mm, fill=omMeth!12] (log) at (11.2, -4.2) {log};
\foreach \i in {0, 1, 2} \draw[<->] (c\i.east) -- (l.west);
\draw[->] (l) -- node[above=2pt, font=\scriptsize] {accept} (q);
\foreach \i in {0, 1, 2, 3} \draw[->] (q.east) -- (w\i.west);
\draw[->, omProp] (w0.east) -- (s0.west); \draw[->, omProp] (w1.east) -- (s2.west); \draw[->, omProp] (w3.east) -- (s3.west);
\draw[->, omMeth] (s3.south) -- (log.north);
\draw[->, omMeth] (s2.east) to[bend left=50] (log.east);
\end{tikzpicture}
\par\smallskip{\small The server's architecture at the end of milestone 4. The listener accepts connections and puts them
in a queue; a fixed set of workers take them, one each, and answer the requests. A request \omterm{def:b2:threads-and-synchronisation:lock}{locks} the one shard that holds
its key; a change is appended to the log under that shard's \omterm{def:b2:threads-and-synchronisation:lock}{lock}.}
\end{omfigure}

\section{Milestone 1: the protocol}\label{sec:b2:project-key-value-server:protocol}

A request is one line of text; the parser turns it into a value of an enum, or into the reason it is not one. Keys and
values are words without spaces, so the arguments are the words of the line.

\omcode{code/bachelor-2/31-project-key-value-server/listings/protocol.rs}{22}{43}{The parser: one line to a
\code{Command}, or the reason in an \code{Err}}\label{lst:b2:project-key-value-server:parse}

The match on the pair (first word, list of arguments) lists the valid forms first, then one error per command with the
wrong number of arguments; the last arm catches every unknown word, lower-case ones included. The enums \code{Command}
and \code{Reply} implement \code{Display}, so that a reply is printed by \code{writeln!} as it is, and a command is written back
as the line it came from: the log of milestone 4 will store commands that way. The unit tests check every command and
eleven malformed lines, one per error reason, and that printing a thousand random commands and parsing them gives the
same commands back.

\section{Milestone 2: one client at a time}\label{sec:b2:project-key-value-server:single}

A connection is a stream of bytes. The server reads it through a buffer, up to the next newline, so that a request
split over two TCP segments, or two requests in one segment, are both handled: the framing is the newline. A line
longer than the limit, or a stream that ends in the middle of a line, cannot be resynchronised, and the connection is
closed after one \code{ERR}.

\omcode{code/bachelor-2/31-project-key-value-server/listings/server.rs}{15}{46}{One connection: read a line, parse it,
apply it, answer it}\label{lst:b2:project-key-value-server:handle}

The replies go through a \code{BufWriter} and are flushed when no request is waiting in the reader's buffer: one
\code{write} for a \omterm{def:b2:link-layer:burst}{burst} of requests, and no reply left waiting while the client waits for it. The simplest server calls
\code{handle} for each accepted connection in turn:

\omcode{code/bachelor-2/31-project-key-value-server/listings/server.rs}{48}{54}{Milestone 2: one connection at a
time}\label{lst:b2:project-key-value-server:one}

The client of \code{client.rs} connects (retrying while the server starts), sends a line and reads one line back. The
demonstration program starts the server of the next milestones in a \omterm{def:b2:threads-and-synchronisation:thread}{thread} of its own \omterm{def:b2:processes-and-system-calls:process}{process}, runs a scripted session
against it, and reads the log back:

\omcode{code/bachelor-2/31-project-key-value-server/listings/demo.rs}{21}{29}{The demonstration: open the log, build
the store from it, serve on port 44310}\label{lst:b2:project-key-value-server:demo}

\omoutput{code/bachelor-2/31-project-key-value-server/outputs/demo.txt}{A scripted session, then the log read back:
four records, one key.}\label{out:b2:project-key-value-server:demo}

Two \code{SET}s, one more that replaces a value, and one \code{DEL} that removes a key: four records, since a \code{DEL}
of a missing key changes nothing and is not logged. Their \omterm{def:b2:project-key-value-server:log}{replay} leaves one key. With one connection at a time, a
second client waits until the first sends \code{QUIT}: a client that connects and stays idle stops the server for
everybody.

\section{Milestone 3: many clients}\label{sec:b2:project-key-value-server:concurrency}

The first remedy gives every connection a \omterm{def:b2:threads-and-synchronisation:thread}{thread} of its own:

\omcode{code/bachelor-2/31-project-key-value-server/listings/server.rs}{56}{64}{A \omterm{def:b2:threads-and-synchronisation:thread}{thread} per
connection}\label{lst:b2:project-key-value-server:per-connection}

The compiler checks the sharing: the closure moves an \code{Arc<Shared>} into the \omterm{def:b2:threads-and-synchronisation:thread}{thread}, and \code{Shared} holds a
\code{Box<dyn Store>} whose \omterm{def:b2:rust-traits-and-generics:trait}{trait} requires \code{Send + Sync}, so a store that could not be shared safely between
\omterm{def:b2:threads-and-synchronisation:thread}{threads} would not compile. A \omterm{def:b2:threads-and-synchronisation:thread}{thread} per connection costs a stack and a creation per client, and nothing bounds the
number of \omterm{def:b2:threads-and-synchronisation:thread}{threads}: ten thousand clients make ten thousand \omterm{def:b2:threads-and-synchronisation:thread}{threads}.

\begin{definition}[Thread pool]\label{def:b2:project-key-value-server:pool}
A \emph{thread pool}\index{thread pool} is a fixed set of \omterm{def:b2:threads-and-synchronisation:thread}{threads}, the \emph{workers}, that take tasks from a shared
queue and run them one after another. Submitting a task puts it in the queue; no \omterm{def:b2:threads-and-synchronisation:thread}{thread} is created per task.
\end{definition}

\omcode{code/bachelor-2/31-project-key-value-server/listings/server.rs}{66}{87}{A pool of workers fed by one
channel}\label{lst:b2:project-key-value-server:pool}

The receiving end of the channel is behind a mutex, since one worker at a time may take from it; the guard is a
temporary of the statement that calls \code{recv}, so the \omterm{def:b2:threads-and-synchronisation:lock}{lock} is released before the connection is handled, not
after. Here a task is a whole connection, so a pool of $w$ workers serves $w$ clients at once and the others wait in the
queue: the pool bounds the resources, at the price of the idle-client problem for the $(w + 1)$-th client. An \omterm{def:b2:sockets:event-loop}{event loop}
(\cref{def:b2:sockets:event-loop}) serves any number of idle clients with one \omterm{def:b2:threads-and-synchronisation:thread}{thread}; the Year 3 volume returns to it.

The store itself must be shared. The first version puts the whole map behind one \code{Mutex}; every request, from
every worker, takes that one \omterm{def:b2:threads-and-synchronisation:lock}{lock}.

\omcode{code/bachelor-2/31-project-key-value-server/listings/store.rs}{12}{44}{The store's \omterm{def:b2:rust-traits-and-generics:trait}{trait}, and a change made
under a \omterm{def:b2:threads-and-synchronisation:lock}{lock} the caller holds}\label{lst:b2:project-key-value-server:store}

\begin{definition}[Sharding, lock striping]\label{def:b2:project-key-value-server:shard}
\emph{Sharding}\index{sharding} splits a map into $N$ parts, the \emph{shards}, and puts each key in the shard numbered
by a hash of the key modulo $N$. With one \omterm{def:b2:threads-and-synchronisation:lock}{lock} per shard, \emph{lock striping}\index{lock striping}, two requests on keys
of different shards proceed in parallel.
\end{definition}

\needspace{16\baselineskip}
\omcode{code/bachelor-2/31-project-key-value-server/listings/store.rs}{109}{122}{The sharded store: a read \omterm{def:b2:threads-and-synchronisation:lock}{lock} for
\code{GET}, a write \omterm{def:b2:threads-and-synchronisation:lock}{lock} for \code{SET} and \code{DEL}}\label{lst:b2:project-key-value-server:sharded}

Each shard is behind a reader--writer \omterm{def:b2:threads-and-synchronisation:lock}{lock}, so that many \code{GET}s on one shard also proceed together. \code{COUNT}
takes the \omterm{def:b2:project-key-value-server:shard}{shards}' read \omterm{def:b2:threads-and-synchronisation:lock}{locks} one at a time; its answer is a sum of counts taken at slightly different instants, which
is the exact count at no single instant if writes happen meanwhile (\cref{exo:b2:project-key-value-server:5} looks at
what an unlocked client sees). Compaction, in the next milestone, holds every shard and then the log: several \omterm{def:b2:threads-and-synchronisation:lock}{locks} at
once.

\begin{proposition}[Lock order]\label{prop:b2:project-key-value-server:lock-order}
Number the \omterm{def:b2:threads-and-synchronisation:lock}{locks}: shard $i$ is \omterm{def:b2:threads-and-synchronisation:lock}{lock} $i$ for $0 \le i < N$, and the log is \omterm{def:b2:threads-and-synchronisation:lock}{lock} $N$. If every operation that holds several
of them acquired them in increasing order, as \code{SET}, \code{DEL}, compaction and the \code{MSET} of
\cref{exo:b2:project-key-value-server:10} do, the server cannot deadlock.
\end{proposition}

\begin{proof}
A \code{SET} or \code{DEL} takes one shard \omterm{def:b2:threads-and-synchronisation:lock}{lock}, then the log: increasing. \code{COUNT} holds one \omterm{def:b2:threads-and-synchronisation:lock}{lock} at a time.
Compaction takes the \omterm{def:b2:project-key-value-server:shard}{shards} in the order $0, 1, \dots, N-1$, then the log. By \cref{prop:b2:threads-and-synchronisation:order},
a set of \omterm{def:b2:threads-and-synchronisation:thread}{threads} waiting for each other would form a circle, in which each \omterm{def:b2:threads-and-synchronisation:thread}{thread} holds a \omterm{def:b2:threads-and-synchronisation:lock}{lock} with a smaller number than
the one it waits for; following the circle, the numbers would increase strictly all the way round and come back to the
first one, which is impossible.
\end{proof}

\Cref{exo:b2:project-key-value-server:6} shows the deadlock that two \code{MSET}s produce when each takes the \omterm{def:b2:threads-and-synchronisation:lock}{locks} in
the order of its own keys.

\section{Milestone 4: persistence}\label{sec:b2:project-key-value-server:persistence}

A map in memory disappears with the \omterm{def:b2:processes-and-system-calls:process}{process}. The server writes every change to a file before it applies it, and
rebuilds the map from that file when it starts.

\begin{definition}[Append-only log, replay, compaction]\label{def:b2:project-key-value-server:log}
An \emph{append-only log}\index{append-only log} is a file to which records are only ever added at the end. To
\emph{replay}\index{log replay} a log of changes is to apply its records in order to an empty state. To
\emph{compact}\index{compaction (log)} it is to replace it by a shorter log whose replay gives the same state.
\end{definition}

A record must be found again even if the file ends in the middle of one, so it is framed by its length, and checked by
its CRC-32:

\begin{omfigure}
\begin{bytefield}[bitwidth=2.1em]{15}
\bitheader{0,3,4,7,8,14} \\
\bitbox{4}{\scriptsize length: 7, little-endian} & \bitbox{4}{\scriptsize CRC-32 of the 7 bytes}
& \bitbox{1}{\scriptsize\code{S}} & \bitbox{1}{\scriptsize\code{E}} & \bitbox{1}{\scriptsize\code{T}}
& \bitbox{1}{\scriptsize\textvisiblespace} & \bitbox{1}{\scriptsize\code{a}} & \bitbox{1}{\scriptsize\textvisiblespace}
& \bitbox{1}{\scriptsize\code{1}}
\end{bytefield}
\par\smallskip{\small The record of \code{SET a 1}, byte by byte: four bytes of length, four of CRC, then the command as
it is printed. A record is 8 bytes longer than its command.}
\end{omfigure}

\omcode{code/bachelor-2/31-project-key-value-server/listings/log.rs}{44}{73}{\omterm{def:b2:project-key-value-server:log}{Replay}: apply each complete record, stop
at the first one that is short, damaged or not a change}\label{lst:b2:project-key-value-server:replay}

The CRC alone would not be enough. A file system may leave zeros at the end of a file after a crash, and four bytes of
length 0 followed by four bytes of CRC 0 form a valid frame, since the CRC-32 of no bytes is 0. The parser catches it: an
empty line is not a command.

\omcode{code/bachelor-2/31-project-key-value-server/listings/log.rs}{81}{101}{Open: \omterm{def:b2:project-key-value-server:log}{replay}, cut off the torn tail, append
from there}\label{lst:b2:project-key-value-server:open}

Cutting the tail off matters: without it, the next record would be appended after the torn bytes, and the following
\omterm{def:b2:project-key-value-server:log}{replay} would stop at the torn bytes and lose the new record too. The policy decides when \code{OK} is sent: after
\code{sync_data}, the record is on the disk; without it, it is in the operating system's memory and lost if the machine
loses power, though not if only the server crashes.

\omcode{code/bachelor-2/31-project-key-value-server/listings/crash_demo.rs}{30}{55}{Three damaged logs, made by
hand}\label{lst:b2:project-key-value-server:crash}

\omoutput{code/bachelor-2/31-project-key-value-server/outputs/crash_demo.txt}{Each restart keeps the longest prefix of
complete, undamaged records.}\label{out:b2:project-key-value-server:crash}

\begin{omfigure}
\begin{tikzpicture}[r/.style={draw, minimum height=6mm, font=\scriptsize\ttfamily, inner sep=1pt},
  >={Stealth[length=1.6mm]}]
\node[font=\small, anchor=east] at (-0.2, 0) {before the crash};
\node[r, minimum width=17mm, fill=omDef!10] (a) at (0.85, 0) {SET a 1};
\node[r, minimum width=17mm, fill=omDef!10] (b) at (2.55, 0) {SET b 2};
\node[r, minimum width=17mm, fill=omDef!10] (c) at (4.25, 0) {SET a 3};
\node[r, minimum width=12mm, fill=omThm!15, dashed] (t) at (5.7, 0) {SET c};
\node[font=\scriptsize, text=omThm, anchor=west] at (6.4, 0) {11 of 15 bytes written};
\node[font=\small, anchor=east] at (-0.2, -1.3) {after the restart};
\node[r, minimum width=17mm, fill=omDef!10] at (0.85, -1.3) {SET a 1};
\node[r, minimum width=17mm, fill=omDef!10] at (2.55, -1.3) {SET b 2};
\node[r, minimum width=17mm, fill=omDef!10] (c2) at (4.25, -1.3) {SET a 3};
\node[r, minimum width=17mm, fill=omMeth!15] at (5.95, -1.3) {SET c 5};
\node[font=\scriptsize, anchor=west] at (6.9, -1.3) {appended after the cut};
\draw[omThm, thick] (5.1, 0.55) -- (5.1, -1.85);
\node[font=\scriptsize, text=omThm] at (5.1, 0.8) {cut};
\node[font=\small, anchor=west] at (0, -2.4) {replay: $a = 3$, $b = 2$; then $c = 5$; the client never received \code{OK} for
\code{SET c 4}};
\end{tikzpicture}
\par\smallskip{\small A crash in the middle of an append. The restart \omterm{def:b2:project-key-value-server:log}{replays} the three complete records, cuts the file
after them, and the next change is appended there.}
\end{omfigure}

\begin{proposition}[Replay]\label{prop:b2:project-key-value-server:replay}
Let a log hold the records $r_1, \dots, r_n$ followed by any bytes that do not begin with a complete, undamaged record.
Its \omterm{def:b2:project-key-value-server:log}{replay} is the map obtained by applying $r_1, \dots, r_n$ in order to the empty map, which is the map the server held
after its $n$-th change.
\end{proposition}

\begin{proof}
\Cref{exo:b2:project-key-value-server:8}, by induction on $n$.
\end{proof}

The proposition assumes that the damage is at the end. A flipped bit in the middle of the log, as in the third case of
\cref{out:b2:project-key-value-server:crash}, also stops the \omterm{def:b2:project-key-value-server:log}{replay}, and the intact records after it are dropped with it:
a crash cannot cause that, a failing disk can, and the honest behaviour is to report it. The kill test of the next
milestone shows why a crash damages only the end: the server appends each record with one \code{write}, and a \omterm{def:b2:processes-and-system-calls:process}{process}
killed by a signal never leaves a write half done; only the machine stopping, between the operating system's memory and
the disk, can.

The log grows with every change, even when the map does not: a thousand \code{SET}s of one key leave a thousand
records for one value. Compaction rewrites it:

\omcode{code/bachelor-2/31-project-key-value-server/listings/log.rs}{103}{123}{Compaction: a new file, made durable,
renamed over the old one}\label{lst:b2:project-key-value-server:compact}

\begin{proposition}[Compaction]\label{prop:b2:project-key-value-server:compaction}
If no change happens during the rewrite, the compacted \omterm{def:b2:project-key-value-server:log}{log replays} to the same map as the log it replaces.
\end{proposition}

\begin{proof}
The compacted log holds one \code{SET k v} for each pair of the map $M$, in sorted order, and no \code{DEL}. Replaying it
inserts each pair once into an empty map, and no two records have the same key, so the order does not matter and the
result is $M$. By \cref{prop:b2:project-key-value-server:replay}, the old \omterm{def:b2:project-key-value-server:log}{log replays} to the map held after its last
change, which is $M$, since no change happens during the rewrite: the store's \code{compact} holds every shard's \omterm{def:b2:threads-and-synchronisation:lock}{lock}
until the new log is in place.
\end{proof}

The rename is what makes compaction safe against a crash: \code{rename} replaces the old name by the new file at once,
so a crash leaves either the old log or the new one, never half of each; syncing the directory makes the rename itself
durable.

\section{Milestone 5: the load test}\label{sec:b2:project-key-value-server:load}

\begin{definition}[Load test, latency percentile]\label{def:b2:project-key-value-server:load-test}
A \emph{load test}\index{load test} runs a fixed, reproducible workload of requests against a system and measures its
\emph{throughput}, the requests completed per second, and the latency of each request. The \emph{$p$-th latency
percentile}\index{latency percentile} is the latency below which a fraction $p$ of the requests fall: $p_{50}$ is the
median, $p_{99}$ the latency that one request in a hundred exceeds.
\end{definition}

The workload is seeded: each client draws its requests from a generator started at its own number, nine \code{GET}s for
one \code{SET} over a thousand keys. It is run twice, once against the store directly from \omterm{def:b2:threads-and-synchronisation:thread}{threads} of the same \omterm{def:b2:processes-and-system-calls:process}{process},
to measure the \omterm{def:b2:threads-and-synchronisation:lock}{locks} alone, and once through TCP:

\omcode{code/bachelor-2/31-project-key-value-server/listings/load.rs}{66}{88}{One run over TCP: each client times every
request; the percentiles come from the sorted times}\label{lst:b2:project-key-value-server:load}

\omoutput{code/bachelor-2/31-project-key-value-server/outputs/load.txt}{One run on the author's laptop: in process, over
TCP, and the cost of syncing the log.}\label{out:b2:project-key-value-server:load}

\begin{omfigure}
\begin{tikzpicture}
\begin{axis}[width=0.62\linewidth, height=5cm, xlabel={threads}, ylabel={million operations per second},
  xtick={1,2,4,8}, xmode=log, log basis x=2, xticklabels={1,2,4,8}, ymin=0, legend pos=north west,
  legend style={font=\small}]
\addplot[mark=*, omThm] table[x=threads, y=one] {code/bachelor-2/31-project-key-value-server/figures/load.dat};
\addplot[mark=square*, omDef] table[x=threads, y=sharded]
  {code/bachelor-2/31-project-key-value-server/figures/load.dat};
\legend{one lock, 16 shards}
\end{axis}
\end{tikzpicture}
\par\smallskip{\small The in-process part of \cref{out:b2:project-key-value-server:load}, measured on a laptop. With one
\omterm{def:b2:threads-and-synchronisation:lock}{lock}, more \omterm{def:b2:threads-and-synchronisation:thread}{threads} do less work in all: they spend their time handing the \omterm{def:b2:threads-and-synchronisation:lock}{lock} to each other. With 16 \omterm{def:b2:project-key-value-server:shard}{shards}, the
throughput grows with the \omterm{def:b2:threads-and-synchronisation:thread}{threads}.}
\end{omfigure}

Three things stand out. In process, one \omterm{def:b2:threads-and-synchronisation:lock}{lock} gets \emph{slower} as \omterm{def:b2:threads-and-synchronisation:thread}{threads} are added, from about 4.5 million operations
per second with one \omterm{def:b2:threads-and-synchronisation:thread}{thread} to about 1 million with eight, while the sharded store reaches about 7 million; with a single
\omterm{def:b2:threads-and-synchronisation:thread}{thread} the \omterm{def:b2:project-key-value-server:shard}{shards} cost something, a hash and a reader--writer \omterm{def:b2:threads-and-synchronisation:lock}{lock} instead of a plain mutex. Over TCP, the difference
nearly vanishes: a request costs on the order of 100 microseconds of \omterm{def:b2:processes-and-system-calls:syscall}{system calls} and loopback, and a \omterm{def:b2:threads-and-synchronisation:lock}{lock} held for a
fraction of a microsecond is not the bottleneck; the throughput grows with the clients because each waits for its own
round trip. And the policy of the log dominates everything when it syncs: a \code{sync_data} per \code{SET} costs
milliseconds, and the server writes some tens of times fewer \code{SET}s than when it leaves the records to the operating
system. Measurements on another machine give other numbers; these three shapes are what to look for.

\begin{inpractice}
A server that must sync can batch the syncs, a technique called \emph{group commit}: the changes that arrive within a
millisecond share one \code{sync}, and every client in the group receives \code{OK} after it. The latency of a \code{SET} grows a little, and the
throughput comes back. Clients also send several requests before reading the replies
(\cref{exo:b2:project-key-value-server:12}).
\end{inpractice}

\omimg{book4/ai/b2-kv-server-rack.jpg}{0.55\linewidth}{A small machine room. A server like this chapter's runs on one
of the machines in the rack; the network cables bring the clients' requests, and the disks keep its log.}

\section{Milestone 6: the test suite, and what is left}\label{sec:b2:project-key-value-server:tests}

The tests form a Cargo crate that includes the chapter's modules by their paths, and a small program that starts the
server on a port and a log given on its command line.

\begin{center}
\small
\begin{tabular}{p{0.3\linewidth}p{0.6\linewidth}}
\hline
test & what it checks \\
\hline
protocol & every command; eleven malformed lines, one per reason; printing then parsing gives the command back \\
stores & both stores agree with a plain map on 3,000 random requests; 8 \omterm{def:b2:threads-and-synchronisation:thread}{threads} of 1,000 \code{SET}s lose none \\
\omterm{def:b2:project-key-value-server:log}{replay} & a log cut at every byte \omterm{def:b2:project-key-value-server:log}{replays} to the map after its last complete record
(\cref{prop:b2:project-key-value-server:replay}); a damaged byte stops the \omterm{def:b2:project-key-value-server:log}{replay} at its record \\
torn tail & \code{open} cuts it and appends after the last good record \\
compaction & the compacted \omterm{def:b2:project-key-value-server:log}{log replays} to the same map, is shorter, and takes appends
(\cref{prop:b2:project-key-value-server:compaction}) \\
\code{MSET} & \omterm{def:b2:threads-and-synchronisation:atomic}{atomic} under 4 writers and 4 readers (\cref{exo:b2:project-key-value-server:10}) \\
end to end & a session on port 44311, a request in two writes, a restart that keeps the data \\
kill & the server killed while a client writes: every acknowledged \code{SET} is replayed \\
\hline
\end{tabular}
\end{center}

The kill test crashes the server on purpose: it starts the server program, lets a client send
\code{SET}s as fast as the syncs allow, and kills the server after 300 milliseconds with a signal that no program can
catch. It then \omterm{def:b2:project-key-value-server:log}{replays} the log and checks that it holds exactly the acknowledged \code{SET}s, plus at most the one in
flight, which was written but whose \code{OK} never left.

What the server does not do: it runs on one machine, so a dead disk or a dead machine loses the data or the service; it
has no authentication, which is acceptable only because it listens on the \omterm{def:b2:sockets:loopback}{loopback address}; it does not expire keys
(\cref{exo:b2:project-key-value-server:7}) or take snapshots (\cref{exo:b2:project-key-value-server:9}); and a
\code{COUNT} is not a snapshot of the map. Replicating the log to other machines, and agreeing on its order when one of
them fails, is the project of the Year 3 volume.

\section{Exercises}\label{sec:b2:project-key-value-server:exercises}

\begin{exercise}[$\star$]\label{exo:b2:project-key-value-server:1}
Predict the server's reply to each line of the session in \code{exo_session.rs}: \code{SET a 1}, \code{SET b 2},
\code{SET a 3}, \code{GET a}, \code{DEL b}, \code{DEL b}, \code{GET b}, \code{COUNT}, \code{SET c 4 5}, \code{SET c 4}
written with two spaces before and after the key, \code{GET c}, \code{DEL}, \code{COUNT}. Then run it.
\end{exercise}

\begin{exercise}[$\star$]\label{exo:b2:project-key-value-server:2}
The server of \code{exo_one_read.rs} treats the bytes returned by one \code{read} as one request. Run it; explain each
line it prints, and what \cref{lst:b2:project-key-value-server:handle} does instead.
\end{exercise}

\begin{exercise}[$\star$]\label{exo:b2:project-key-value-server:3}
A visit counter is incremented by \code{GET} and then \code{SET} of the value plus one. Explain the result of the first
part of \code{exo_incr.rs}, then write an \code{INCR k} that is \omterm{def:b2:threads-and-synchronisation:atomic}{atomic}, and check it with eight clients.
\end{exercise}

\begin{exercise}[$\star$]\label{exo:b2:project-key-value-server:4}
Keys have 4 characters and values 6. Give the size of the log after $n$ \code{SET}s spread evenly over $k$ keys, before
and after compaction, as formulas in $n$ and $k$; check them with \code{exo_log_size.rs}.
\end{exercise}

\begin{exercise}[$\star\star$]\label{exo:b2:project-key-value-server:5}
Two clients run without any \omterm{def:b2:threads-and-synchronisation:lock}{lock}: each reads $x$ and then writes the value it read plus one, starting from $x = 0$. List
the \omterm{def:b2:threads-and-synchronisation:interleaving}{interleavings} of the four steps, give the final $x$ for each, and check with \code{exo_interleave.rs}. Which outcome
is the \omterm{def:b2:database-design:anomalies}{lost update} of the previous chapter?
\end{exercise}

\begin{exercise}[$\star\star$]\label{exo:b2:project-key-value-server:6}
An \code{MSET a x b y} must change two keys in different \omterm{def:b2:project-key-value-server:shard}{shards} together. A first version \omterm{def:b2:threads-and-synchronisation:lock}{locks} the shard of its first
key, then the shard of its second. Run \code{exo_mset_deadlock.rs}, explain the first two lines with a wait-for graph,
and fix the order.
\end{exercise}

\begin{exercise}[$\star\star$]\label{exo:b2:project-key-value-server:7}
Add \code{SETEX k s v}, which sets a key that disappears $s$ seconds later. Run \code{exo_setex.rs}, which expires a key
when it is read and sweeps the others. What must the log record so that a restart does not bring an expired key back?
\end{exercise}

\begin{exercise}[$\star\star$]\label{exo:b2:project-key-value-server:8}
Prove \cref{prop:b2:project-key-value-server:replay}.
\end{exercise}

\begin{exercise}[$\star\star$]\label{exo:b2:project-key-value-server:9}
Instead of compacting, the server can write a \emph{snapshot}, the whole map as a file, and start a new, empty log.
Run \code{exo_snapshot.rs}. In which order must the snapshot and the new log be made durable, and what does recovery do
if the server crashes while writing the snapshot?
\end{exercise}

\begin{exercise}[$\star\star\star$]\label{exo:b2:project-key-value-server:10}
Write an \omterm{def:b2:threads-and-synchronisation:atomic}{atomic} \code{MSET} and \code{MGET} for the sharded store, taking the shard \omterm{def:b2:threads-and-synchronisation:lock}{locks} in increasing order, and test
them under load: four writers set two keys to the same value, half of them naming the keys in the other order, and four
readers check that they never see two different values. Use \code{exo_mset.rs} and its module \code{mset.rs}.
\end{exercise}

\begin{exercise}[$\star\star\star$]\label{exo:b2:project-key-value-server:11}
Model the server as follows: a request costs $s$ microseconds in all, of which a time $h$ is spent holding the
\omterm{def:b2:threads-and-synchronisation:lock}{lock} of its shard; the store has $N$ \omterm{def:b2:project-key-value-server:shard}{shards} and receives requests from $c$ clients that each wait for the reply before
sending the next. Give the throughput when the \omterm{def:b2:threads-and-synchronisation:lock}{lock} is not the bottleneck, and the throughput that one \omterm{def:b2:threads-and-synchronisation:lock}{lock} caps, and
compare with \cref{out:b2:project-key-value-server:load}.
\end{exercise}

\begin{exercise}[$\star\star\star$]\label{exo:b2:project-key-value-server:12}
Pipelining: a client sends a batch of $b$ requests in one write, then reads the $b$ replies. Measure the throughput for
$b = 1, 10, 100, 1000$ with \code{exo_pipeline.rs}, and explain the result. Why did \cref{lst:b2:project-key-value-server:handle}
need nothing new to support it?
\end{exercise}
